% GENERATED FILE. Edit canonical lessons, then run npm run generate:books.
% canonical-source-hash: fnv1a64:5e3e333dc6c0505a
% canonical-lessons: PT-C302-sismo, PT-C302-terramoto, PT-C302-tremor, PT-C302-ruir, PT-C302-escombros

\chapter{A terra tremeu}
\label{ch:pt-a-terra-tremeu}
\hlchaptermodality{302}

\begin{chapteropening}
\textbf{By the end of this chapter:} \emph{I can talk about earthquakes, tremors, buildings that collapsed and the rubble.}

Complete the last lesson of chapter 302: I can talk about earthquakes, tremors, buildings that collapsed and the rubble.
\end{chapteropening}

\section[o sismo]{o sismo — the quake (formal)}

\label{lesson:PT-C302-sismo}



\begin{quote}
Before the new one: say the Portuguese for the fire extinguisher, then the Portuguese for the burn.
\end{quote}

\subsection*{You'll want to know: o sismo}
\textbf{o sismo} — \textquotedblleft{}the quake (formal)\textquotedblright{}.

\begin{quote}
O sismo foi sentido em todo o país.
\end{quote}

\begin{quote}
\emph{The quake was felt all over the country.}
\end{quote}

A news word in Portugal.

\begin{cousinweb}
\emph{sismo} was borrowed from Greek \emph{seismós} \textquotedblleft{}a shaking, an earthquake\textquotedblright{} (EtymDB), from \emph{seíein} \textquotedblleft{}to shake\textquotedblright{}.
\end{cousinweb}

\begin{grammarlens}[title={What you've built}]
One more word for telling, explaining and describing at B1.
\end{grammarlens}

\subsection*{Your turn}
\begin{itemize}
\raggedright
  \item \emph{Say it:} \emph{o sismo}
  \item \emph{Say it:} \emph{o sismo}, once more
  \item \emph{Say it:} a sentence of your own with \emph{o sismo}
  \item \emph{Recall:} say the Portuguese for the fire extinguisher, then the Portuguese for the burn, then say \emph{o sismo} again
\end{itemize}

\begin{tcolorbox}[breakable,colback=teal!4,colframe=teal!35!black,title={Before you move on}]
What does o sismo mean? (\textquotedblleft{}The quake (formal)\textquotedblright{}.) Say it once more.
\end{tcolorbox}
\section[o terramoto]{o terramoto — the earthquake}

\label{lesson:PT-C302-terramoto}



\begin{quote}
Before the new one: say the Portuguese for the burn, then the Portuguese for the quake.
\end{quote}

\subsection*{You'll want to know: o terramoto}
\textbf{o terramoto} — \textquotedblleft{}the earthquake\textquotedblright{}.

\begin{quote}
Em 1755, um terramoto destruiu a cidade.
\end{quote}

\begin{quote}
\emph{In 1755 an earthquake destroyed the city.}
\end{quote}

In Brazil: \emph{terremoto}.

\begin{cousinweb}
\emph{terramoto} goes back to Latin \emph{terrae mōtus} \textquotedblleft{}a movement of the earth\textquotedblright{}, from \emph{terra} \textquotedblleft{}earth\textquotedblright{} and \emph{movēre} \textquotedblleft{}to move\textquotedblright{} (Lewis \& Short).
\end{cousinweb}

\begin{grammarlens}[title={What you've built}]
One more word for telling, explaining and describing at B1.
\end{grammarlens}

\subsection*{Your turn}
\begin{itemize}
\raggedright
  \item \emph{Say it:} \emph{o terramoto}
  \item \emph{Say it:} \emph{o terramoto}, once more
  \item \emph{Say it:} a sentence of your own with \emph{o terramoto}
  \item \emph{Recall:} say the Portuguese for the burn, then the Portuguese for the quake, then say \emph{o terramoto} again
\end{itemize}

\begin{tcolorbox}[breakable,colback=teal!4,colframe=teal!35!black,title={Before you move on}]
What does o terramoto mean? (\textquotedblleft{}The earthquake\textquotedblright{}.) Say it once more.
\end{tcolorbox}
\section[o tremor]{o tremor — the tremor, the trembling}

\label{lesson:PT-C302-tremor}



\begin{quote}
Before the new one: say the Portuguese for the quake, then the Portuguese for the earthquake.
\end{quote}

\subsection*{You'll want to know: o tremor}
\textbf{o tremor} — \textquotedblleft{}the tremor, the trembling\textquotedblright{}.

\begin{quote}
Sentimos um tremor às três da manhã.
\end{quote}

\begin{quote}
\emph{We felt a tremor at three in the morning.}
\end{quote}

\emph{Tremor de terra} — earth tremor.

\begin{cousinweb}
\emph{tremor} comes from Latin \emph{tremor} \textquotedblleft{}a trembling\textquotedblright{} (EtymDB), from \emph{tremere} \textquotedblleft{}to tremble\textquotedblright{}.
\end{cousinweb}

\begin{grammarlens}[title={What you've built}]
One more word for telling, explaining and describing at B1.
\end{grammarlens}

\subsection*{Your turn}
\begin{itemize}
\raggedright
  \item \emph{Say it:} \emph{o tremor}
  \item \emph{Say it:} \emph{o tremor}, once more
  \item \emph{Say it:} a sentence of your own with \emph{o tremor}
  \item \emph{Recall:} say the Portuguese for the quake, then the Portuguese for the earthquake, then say \emph{o tremor} again
\end{itemize}

\begin{tcolorbox}[breakable,colback=teal!4,colframe=teal!35!black,title={Before you move on}]
What does o tremor mean? (\textquotedblleft{}The tremor, the trembling\textquotedblright{}.) Say it once more.
\end{tcolorbox}
\section[ruir]{ruir — to collapse, to fall down}

\label{lesson:PT-C302-ruir}



\begin{quote}
Before the new one: say the Portuguese for the earthquake, then the Portuguese for the tremor.
\end{quote}

\subsection*{You'll want to know: ruir}
\textbf{ruir} — \textquotedblleft{}to collapse, to fall down\textquotedblright{}.

\begin{quote}
O velho prédio ruiu durante a noite.
\end{quote}

\begin{quote}
\emph{The old building collapsed during the night.}
\end{quote}

Mostly in the third person: \emph{ruiu, ruíram}.

\begin{cousinweb}
\emph{ruir} comes from Latin \emph{ruere} \textquotedblleft{}to rush down, to fall\textquotedblright{} (EtymDB).
\end{cousinweb}

\begin{grammarlens}[title={What you've built}]
One more word for telling, explaining and describing at B1.
\end{grammarlens}

\subsection*{Your turn}
\begin{itemize}
\raggedright
  \item \emph{Say it:} \emph{ruir}
  \item \emph{Say it:} \emph{ruir}, once more
  \item \emph{Say it:} a sentence of your own with \emph{ruir}
  \item \emph{Recall:} say the Portuguese for the earthquake, then the Portuguese for the tremor, then say \emph{ruir} again
\end{itemize}

\begin{tcolorbox}[breakable,colback=teal!4,colframe=teal!35!black,title={Before you move on}]
What does ruir mean? (\textquotedblleft{}To collapse, to fall down\textquotedblright{}.) Say it once more.
\end{tcolorbox}
\section[os escombros]{os escombros — the rubble, the ruins}

\label{lesson:PT-C302-escombros}



\begin{quote}
Before the new one: say the Portuguese for the tremor, then the Portuguese for to collapse.
\end{quote}

\subsection*{You'll want to know: os escombros}
\textbf{os escombros} — \textquotedblleft{}the rubble, the ruins\textquotedblright{}.

\begin{quote}
Os bombeiros procuram pessoas nos escombros.
\end{quote}

\begin{quote}
\emph{The firefighters are searching the rubble for people.}
\end{quote}

Almost always plural.

\begin{cousinweb}
The origin of \emph{escombros} is not settled; it sits beside Spanish \emph{escombro}, and EtymDB has no entry.
\end{cousinweb}

\begin{grammarlens}[title={What you've built}]
One more word for telling, explaining and describing at B1.
\end{grammarlens}

\subsection*{Your turn}
\begin{itemize}
\raggedright
  \item \emph{Say it:} \emph{os escombros}
  \item \emph{Say it:} \emph{os escombros}, once more
  \item \emph{Say it:} a sentence of your own with \emph{os escombros}
  \item \emph{Recall:} say the Portuguese for the tremor, then the Portuguese for to collapse, then say \emph{os escombros} again
\end{itemize}

\begin{tcolorbox}[breakable,colback=teal!4,colframe=teal!35!black,title={Before you move on}]
What does os escombros mean? (\textquotedblleft{}The rubble, the ruins\textquotedblright{}.) Say it once more.
\end{tcolorbox}
